\subsection{Auxiliary Regressions and Moments}\label{sec:appendix_moments}
In this section we list the parameters of the auxiliary models and the additional moments we match in estimation. The standard set of controls in the regressions is: age, gender, very-low-SES index ({\itshape alumno prioritario}), dummy for whether the student ever failed a grade, school-track type, baseline SIMCE score.
\begin{enumerate}[leftmargin=5.5mm]
    \item {\itshape Treatment Effect Regressions:}
    
    \begin{itemize}[leftmargin=3.5mm] \setlength\itemsep{0.0em}
        \item All parameters, including the constant, of a regression of achievement on treatment, the standard controls, and average GPA in $9^{th}$ and $10^{th}$ grade (9).
        \item Coefficient on treatment of a regression of hours of study on treatment and the standard controls (1).
        \item Coefficient on treatment of a regression of hours of study on treatment and the standard controls for the sample of students who report, at baseline, no intention to attend college (1).
        \item Coefficient on treatment of a regression of college enrollment on treatment and the standard controls (1).
        \item Coefficient on treatment of a regression of taking the entrance exam on treatment and the standard controls (1).
    \end{itemize}
    
    
    \item {\itshape Descriptive Regressions:}
    \begin{itemize}[leftmargin=3.5mm]\setlength\itemsep{0.0em}
     \item Constant and coefficient of regression of hours of study on dummy for whether student has no intention to stay in school beyond high school (2).
     \item Coefficient on $10^{th}$ grade GPA of regression of $12^{th}$ grade GPA on $10^{th}$ grade GPA (1).
     \item Coefficient on baseline SIMCE score of regression of entrance exam score on baseline SIMCE score (1).
     \item Coefficients on whether the student participated in the survey and on the average between $9^{th}$ and $10^{th}$ grade GPA in a regression of whether a student takes the entrance exam on these variables and on the standard controls (2).
     \item Coefficient on the average between $9^{th}$ and $10^{th}$ grade GPA in a regression of study hours on this variable and on the standard controls (1).
    \end{itemize}
    
    \item {\itshape Descriptive Statistics:}
    \begin{itemize}[leftmargin=3.5mm]\setlength\itemsep{0.0em}
        \item Mean and variance of hours of study (2).
        \item Fraction of students admitted to college by treatment group and baseline achievement, i.e., above or below median SIMCE score (4).
        \item Correlation between regular admissions and PACE admissions for treated students (1).
        \item Fraction taking entrance exam by treatment group (2).
        \item Mean and variance of entrance exam score by treatment group (4).
        \item Fraction of students who enroll in college by treatment group and baseline achievement, i.e., above or below median SIMCE score (4).
        \item Fraction of students enrolled in college by very-low-SES status, i.e., {\itshape alumno prioritario} categorization (2).
        \item Mean and variance of GPA in the control group (2).
        \item All pairwise correlations between the expected score on the PSU, enrollment, and the actual score on the PSU (3).
        \item Mean and variance of perceived returns to effort in GPA production and in PSU production (4).
        \item Correlation between taking the entrance exam and enrollment in the control group (1).
        \item Correlation between study hours and enrollment in the control group (1).
        \item Correlation between study hours and admissions in the control group (1).
        \item Correlation between taking the entrance exam and perceived distance from the within-school cutoff in the treatment group (1).
        \item Correlation between taking the entrance exam and expected PSU score in the control group (1).
        \item Unexplained variation in achievement and GPA after controlling for all initial conditions in the model affecting these outcomes. Specifically, variance of the residuals from regressions of achievement and of GPA on treatment, GPA in $9^{th}$ grade and average GPA between $9^{th}$ and $10^{th}$ grade, a dummy for whether a student reported at baseline to not being interested in attending college, perceived within-school cutoff, and the standard controls (2).
        \item Fractions enrolling through the regular and through the PACE channel for those admitted through both channels (2).
        \item Selectivity of the regular and of the PACE admissions for those admitted through both channels (2).
        \item Mean and variance of expected GPA and PSU score (4).
    \end{itemize}

    


\end{enumerate}